\documentclass[10pt,a4paper]{amsart}
\usepackage{amsmath,amssymb,amsthm,graphicx}
% Retained minimal preamble for the Stacks Project chapter flat.tex (More on Flatness).
% Verbatim excerpt of the chapter preamble: Xy-pic setup (native preamble lines 14-18) and the native
% theorem declarations with the mathematical macros (native preamble lines 154-217). Copied mechanically;
% nothing added, removed or redefined. The wrapper class and the amsmath/amssymb/amsthm/graphicx base
% packages are supplied by the assembler, so the class/externaldocument part of the native preamble is omitted.
% For commutative diagrams we use Xy-pic
\usepackage[all]{xy}

% We use 2cell for 2-commutative diagrams.
\xyoption{2cell}
% Theorem environments.
%
\theoremstyle{plain}
\newtheorem{theorem}[subsection]{Theorem}
\newtheorem{proposition}[subsection]{Proposition}
\newtheorem{lemma}[subsection]{Lemma}

\theoremstyle{definition}
\newtheorem{definition}[subsection]{Definition}
\newtheorem{example}[subsection]{Example}
\newtheorem{exercise}[subsection]{Exercise}
\newtheorem{situation}[subsection]{Situation}

\theoremstyle{remark}
\newtheorem{remark}[subsection]{Remark}
\newtheorem{remarks}[subsection]{Remarks}

\numberwithin{equation}{subsection}

% Macros
%
\def\lim{\mathop{\mathrm{lim}}\nolimits}
\def\colim{\mathop{\mathrm{colim}}\nolimits}
\def\Spec{\mathop{\mathrm{Spec}}}
\def\Hom{\mathop{\mathrm{Hom}}\nolimits}
\def\Ext{\mathop{\mathrm{Ext}}\nolimits}
\def\SheafHom{\mathop{\mathcal{H}\!\mathit{om}}\nolimits}
\def\SheafExt{\mathop{\mathcal{E}\!\mathit{xt}}\nolimits}
\def\Sch{\mathit{Sch}}
\def\Mor{\mathop{\mathrm{Mor}}\nolimits}
\def\Ob{\mathop{\mathrm{Ob}}\nolimits}
\def\Sh{\mathop{\mathit{Sh}}\nolimits}
\def\NL{\mathop{N\!L}\nolimits}
\def\CH{\mathop{\mathrm{CH}}\nolimits}
\def\proetale{{pro\text{-}\acute{e}tale}}
\def\etale{{\acute{e}tale}}
\def\QCoh{\mathit{QCoh}}
\def\Ker{\mathop{\mathrm{Ker}}}
\def\Im{\mathop{\mathrm{Im}}}
\def\Coker{\mathop{\mathrm{Coker}}}
\def\Coim{\mathop{\mathrm{Coim}}}

% Boxtimes
%
\DeclareMathSymbol{\boxtimes}{\mathbin}{AMSa}{"02}

%
% Macros for moduli stacks/spaces
%
\def\QCohstack{\mathcal{QC}\!\mathit{oh}}
\def\Cohstack{\mathcal{C}\!\mathit{oh}}
\def\Spacesstack{\mathcal{S}\!\mathit{paces}}
\def\Quotfunctor{\mathrm{Quot}}
\def\Hilbfunctor{\mathrm{Hilb}}
\def\Curvesstack{\mathcal{C}\!\mathit{urves}}
\def\Polarizedstack{\mathcal{P}\!\mathit{olarized}}
\def\Complexesstack{\mathcal{C}\!\mathit{omplexes}}
% \Pic is the operator that assigns to X its picard group, usage \Pic(X)
% \Picardstack_{X/B} denotes the Picard stack of X over B
% \Picardfunctor_{X/B} denotes the Picard functor of X over B
\def\Pic{\mathop{\mathrm{Pic}}\nolimits}
\def\Picardstack{\mathcal{P}\!\mathit{ic}}
\def\Picardfunctor{\mathrm{Pic}}
\def\Deformationcategory{\mathcal{D}\!\mathit{ef}}

% --- appended 2026-10-12 (Chief): the project own definitions of the reference/slogan/history
% display environments, copied mechanically from the shared Stacks preamble lines 8-12
% (\usepackage{verbatim} plus the three \newenvironment lines, which the chapter preamble keeps
% verbatim from the project preamble). Nothing added, removed or redefined: these environments
% carry no mathematical content and the project itself typesets them as comments. Without these
% three lines a counted statement whose own text uses them fails with Environment slogan undefined.
\usepackage{verbatim}
\newenvironment{reference}{\comment}{\endcomment}
\newenvironment{slogan}{\comment}{\endcomment}
\newenvironment{history}{\comment}{\endcomment}
\usepackage{cleveref}
\crefname{originalreference}{original reference}{original references}
\Crefname{originalreference}{Original reference}{Original references}
\makeatletter
\newlabel{situation-existence}{{0CTC}{source}{}{sourcepointer-1}{}}
\newlabel{situation-existence@cref}{{[originalreference][0][]0CTC}{[1][0][]source}{}{}{}}
\makeatother
\crefname{originallemma}{original lemma}{original lemmas}
\Crefname{originallemma}{Original lemma}{Original lemmas}
\makeatletter
\newlabel{lemma-compute-what-it-should-be}{{0CTD}{source}{}{sourcepointer-2}{}}
\newlabel{lemma-compute-what-it-should-be@cref}{{[originallemma][0][]0CTD}{[1][0][]source}{}{}{}}
\makeatother
\crefname{originalsection}{original section}{original sections}
\Crefname{originalsection}{Original section}{Original sections}
\makeatletter
\newlabel{cohomology-section-global-RHom}{{0B6A}{source}{}{sourcepointer-3}{}}
\newlabel{cohomology-section-global-RHom@cref}{{[originalsection][0][]0B6A}{[1][0][]source}{}{}{}}
\makeatother
\crefname{originallemma}{original lemma}{original lemmas}
\Crefname{originallemma}{Original lemma}{Original lemmas}
\makeatletter
\newlabel{perfect-lemma-base-change-RHom-perfect}{{0A1J}{source}{}{sourcepointer-4}{}}
\newlabel{perfect-lemma-base-change-RHom-perfect@cref}{{[originallemma][0][]0A1J}{[1][0][]source}{}{}{}}
\makeatother
\crefname{originallemma}{original lemma}{original lemmas}
\Crefname{originallemma}{Original lemma}{Original lemmas}
\makeatletter
\newlabel{more-algebra-lemma-Rlim-perfect-gives-perfect}{{0CQG}{source}{}{sourcepointer-5}{}}
\newlabel{more-algebra-lemma-Rlim-perfect-gives-perfect@cref}{{[originallemma][0][]0CQG}{[1][0][]source}{}{}{}}
\makeatother
\crefname{originallemma}{original lemma}{original lemmas}
\Crefname{originallemma}{Original lemma}{Original lemmas}
\makeatletter
\newlabel{cohomology-lemma-dual-perfect-complex}{{08DQ}{source}{}{sourcepointer-6}{}}
\newlabel{cohomology-lemma-dual-perfect-complex@cref}{{[originallemma][0][]08DQ}{[1][0][]source}{}{}{}}
\makeatother
\title{Global sections and derived Hom against a perfect object in the existence situation are perfect and compute the reductions}
\author{The Stacks Project Authors}
\date{stacks-project:flat.tex:local-snapshot-20261011; GNU FDL 1.2 or later}
\begin{document}
\maketitle
\noindent\textit{Editorial note.} All delivered mathematical text is the author's own, copied from the pinned Stacks Project chapter file flat.tex; no statement or proof step was rewritten, reordered or removed.\par\medskip
\noindent\textit{Editorial note.} The preamble is the chapter's own retained minimal preamble (sources/stacks-flat/retained-preamble.tex), a verbatim excerpt of the native Xy-pic setup and of the native theorem declarations and mathematical macros, followed by the project own definitions of the reference, slogan and history display environments (verbatim package plus three newenvironment lines) that the retained file carries at its end and that the project itself typesets as comments. No macro was invented or redefined, and the wrapper class and the amsmath/amssymb/amsthm/graphicx base packages are supplied by the assembler.\par\medskip
\noindent\textit{Editorial note.} The counted statement and proof contain no cross reference, no citation and no included file, so the delivered body is self contained.\par\medskip
\noindent\textit{Editorial note.} Canonical location of the counted proposition: Stacks tag 0CTE (sources/stacks-flat/pinned-original-source.tex line 8079).\par\medskip
\noindent\textit{Allowed standard prerequisites (editorial).}\begin{itemize}
\item The object K and the pendant lemma. Same chapter, Lemma labelled lemma-compute-what-it-should-be, invoked in the statement. Supplies K and the system F\_n used throughout.
\item Base change for derived Hom of perfect complexes. Derived Categories of Schemes, Lemma labelled perfect-lemma-base-change-RHom-perfect, cited at proof lines 8104-8106. Shows each N\_n is perfect over D(A\_n) with compatible base change.
\item Derived limits of perfect complexes. More on Algebra, Lemma labelled more-algebra-lemma-Rlim-perfect-gives-perfect, cited at proof lines 8110-8114. Makes Rlim N\_n perfect and recovers its base changes.
\item Derived Hom as global sections of the sheaf Hom, and duality with the tensor product formulation. Cohomology, Section labelled cohomology-section-global-RHom, cited at proof lines 8100-8102, and Lemma labelled cohomology-lemma-dual-perfect-complex, cited at proof lines 8123-8125. Identifies the functors and translates statement (1) into statement (2).
\end{itemize}\medskip
\ifdefined\hypertarget\hypertarget{sourcepointer-1}{}\fi
\noindent\textit{Source reference.} Original reference 0CTC: sources/stacks-flat/pinned-original-source.tex, source lines 8036--8036. Same-chapter pointer into the author environment situation at source line 8036; invoked at proof line None of the counted proof. It is retained only as a navigation pointer; every load-bearing step is inside the delivered span.\par\medskip
\ifdefined\hypertarget\hypertarget{sourcepointer-2}{}\fi
\noindent\textit{Source reference.} Original lemma 0CTD: sources/stacks-flat/pinned-original-source.tex, source lines 8062--8062. Same-chapter pointer into the author environment lemma at source line 8062; invoked at proof line None of the counted proof. It is retained only as a navigation pointer; every load-bearing step is inside the delivered span.\par\medskip
\ifdefined\hypertarget\hypertarget{sourcepointer-3}{}\fi
\noindent\textit{Source reference.} Original section 0B6A: sources/stacks-chow/tags.txt, source lines 14063--14063. Cross-chapter standard prerequisite located in the project tag list staged in the local snapshot; invoked at proof line 8103 of the counted proof. It is retained only as a navigation pointer; every load-bearing step is inside the delivered span.\par\medskip
\ifdefined\hypertarget\hypertarget{sourcepointer-4}{}\fi
\noindent\textit{Source reference.} Original lemma 0A1J: sources/stacks-chow/tags.txt, source lines 12672--12672. Cross-chapter standard prerequisite located in the project tag list staged in the local snapshot; invoked at proof line 8105 of the counted proof. It is retained only as a navigation pointer; every load-bearing step is inside the delivered span.\par\medskip
\ifdefined\hypertarget\hypertarget{sourcepointer-5}{}\fi
\noindent\textit{Source reference.} Original lemma 0CQG: sources/stacks-chow/tags.txt, source lines 15961--15961. Cross-chapter standard prerequisite located in the project tag list staged in the local snapshot; invoked at proof line 8110 of the counted proof. It is retained only as a navigation pointer; every load-bearing step is inside the delivered span.\par\medskip
\ifdefined\hypertarget\hypertarget{sourcepointer-6}{}\fi
\noindent\textit{Source reference.} Original lemma 08DQ: sources/stacks-chow/tags.txt, source lines 10645--10645. Cross-chapter standard prerequisite located in the project tag list staged in the local snapshot; invoked at proof line 8124 of the counted proof. It is retained only as a navigation pointer; every load-bearing step is inside the delivered span.\par\medskip
\begin{lemma}
\label{lemma-compute-against-perfect}
In Situation \ref{situation-existence} let $K$ be as in
Lemma \ref{lemma-compute-what-it-should-be}. For any perfect
object $E$ of $D(\mathcal{O}_X)$ we have
\begin{enumerate}
\item $M = R\Gamma(X, K \otimes^\mathbf{L} E)$ is a perfect object of $D(A)$
and there is a canonical isomorphism
$R\Gamma(X_n, \mathcal{F}_n \otimes^\mathbf{L} E|_{X_n}) =
M \otimes_A^\mathbf{L} A_n$
in $D(A_n)$,
\item $N = R\Hom_X(E, K)$ is a perfect object of $D(A)$
and there is a canonical isomorphism
$R\Hom_{X_n}(E|_{X_n}, \mathcal{F}_n) = N \otimes_A^\mathbf{L} A_n$
in $D(A_n)$.
\end{enumerate}
In both statements $E|_{X_n}$ denotes the derived pullback
of $E$ to $X_n$.
\end{lemma}

\begin{proof}
Proof of (2). Write $E_n = E|_{X_n}$ and
$N_n = R\Hom_{X_n}(E_n, \mathcal{F}_n)$.
Recall that $R\Hom_{X_n}(-, -)$ is equal to
$R\Gamma(X_n, R\SheafHom(-, -))$, see
Cohomology, Section \ref{cohomology-section-global-RHom}.
Hence by Derived Categories of Schemes, Lemma
\ref{perfect-lemma-base-change-RHom-perfect}
we see that $N_n$ is a perfect object of $D(A_n)$
whose formation commutes with base change. Thus the maps
$N_n \otimes_{A_n}^\mathbf{L} A_{n - 1} \to N_{n - 1}$
coming from $\varphi_n$ are isomorphisms.
By More on Algebra, Lemma \ref{more-algebra-lemma-Rlim-perfect-gives-perfect}
we find that $R\lim N_n$ is perfect and
that its base change back to $A_n$ recovers $N_n$.
On the other hand, the exact functor
$R\Hom_X(E, -) : D_\QCoh(\mathcal{O}_X) \to D(A)$
of triangulated categories commutes with products
and hence with derived limits, whence
$$
R\Hom_X(E, K) =
R\lim R\Hom_X(E, \mathcal{F}_n) =
R\lim R\Hom_X(E_n, \mathcal{F}_n) =
R\lim N_n
$$
This proves (2). To see that (1) holds, translate it into (2)
using Cohomology, Lemma \ref{cohomology-lemma-dual-perfect-complex}.
\end{proof}

\end{document}
